\section{Implementation details and theoretical correspondence}
\label{sec:implementation-correspondence}
This appendix specifies the observations, representation models, and
training operations summarized in Section~\ref{sec:method}. It also
records their correspondence with the population analysis and distinguishes
the synthetic and real-video implementations.

\subsection{Physical simulator and implemented observations}
\label{subsec:implementation-process}

We distinguish the process that generates a video from the model that extracts
its action signal. Let $\mathbf c_{\mathrm{scene}}$ denote persistent scene properties,
$\mathbf q_t$ a physical state sufficient to evolve the scene, and $u$ an action regime.
Unlike the appearance-only $\mathbf z^{\mathrm c}$ in Appendix~\ref{sec:controlled-transfer},
$\mathbf c_{\mathrm{scene}}$ may include properties that affect physical dynamics.
The physical state includes configurations and velocities of interacting
entities; it is not a diffusion latent or a learned action embedding.
A controlled dynamical model is
\begin{align}
  (\mathbf c_{\mathrm{scene}},\mathbf q_0,u)&\sim p_0,\qquad
       \mathbf u_t^{\mathrm{ctrl}}=u_{\mathrm{ctrl}}(u,t), \label{eq:impl-context}\\
  \mathbf q_t&=F_{\mathrm{phys}}(\mathbf c_{\mathrm{scene}},\mathbf q_{t-1},\mathbf u_t^{\mathrm{ctrl}},\boldsymbol\epsilon_t),
       \label{eq:impl-action}\\
  \mathbf x_t&=g_{\mathrm{render}}(\mathbf c_{\mathrm{scene}},\mathbf q_t,\mathbf n_t,\boldsymbol\nu_t),
       \qquad t=1,\ldots,T. \label{eq:impl-observation}
\end{align}
Here $\mathbf u_t^{\mathrm{ctrl}}$ is a control input, $\mathbf n_t$ denotes time-varying capture conditions,
and $\boldsymbol\epsilon_t,\boldsymbol\nu_t$ denote transition and observation noise, respectively.
The joint law $p_0$ allows scene properties and action choice to be correlated.
These equations do not impose componentwise independent innovations,
a sparse causal graph, or an invertible renderer.

\paragraph{Instantiation in the synthetic generator.}
The implemented simulator samples scene parameters from a fixed seed,
initializes a MuJoCo scene~\citep{todorov2012mujoco}, and supplies an action-dependent tool position
at each frame. Physics substeps evolve the objects before rendering.
Thus $F_{\mathrm{phys}}$ is the simulator transition, $u_{\mathrm{ctrl}}$
is the prescribed controller, and $g_{\mathrm{render}}$ is the renderer.
Conditional on the scene and control schedule, this implementation introduces
no explicit stochastic transition or observation noise.
Its camera and lighting remain fixed within a clip and can be absorbed into
$\mathbf c_{\mathrm{scene}}$. Randomness across clips comes from scene construction,
rather than independent noise added to every action coordinate.

Writing $\mathcal V_{\mathrm{sim}}(\mathbf c_{\mathrm{scene}},\mathbf q_0,u)$ for the resulting
rollout, synthetic editing pairs satisfy
\begin{equation}
 \begin{aligned}
  V_{\mathrm s}&=\mathcal V_{\mathrm{sim}}(\mathbf c_{\mathrm{scene}},\mathbf q_0,u_{\mathrm s}),\\
  V^{\mathrm{cf}}&=\mathcal V_{\mathrm{sim}}(\mathbf c_{\mathrm{scene}},\mathbf q_0,u^\star).
 \end{aligned}
 \label{eq:impl-counterfactual}
\end{equation}
The implementation reuses the scene and initial state and changes the
control schedule; interactions and object trajectories are then recomputed.
The instruction names the requested action. No learned decomposition of
text into semantic atoms is required by this construction.
For real videos, Eqs.~\ref{eq:impl-context}--\ref{eq:impl-observation} are a
modeling abstraction. Collected source--target clips need not share the
initial state or capture conditions and are not automatically samples of
the strict coupling in Eq.~\ref{eq:impl-counterfactual}.

\paragraph{Physical state versus action descriptor.}
The synthetic representation experiment constructs a supervised descriptor
\begin{equation}
  \mathbf d_{1:L_{\mathrm{syn}}}=\mathcal T(\mathbf q_{0:T},\mathbf u^{\mathrm{ctrl}}_{1:T}),
  \qquad \mathbf d_t\in\R^{48},\quad L_{\mathrm{syn}}=32.
  \label{eq:impl-synthetic-descriptor}
\end{equation}
The fixed transformation $\mathcal T$ resamples simulator traces and records
displacements, finite-difference velocities and accelerations, relative
geometry, contacts, and interaction phases and order. Positions are rebased
to the first phase. Here $t=1,\ldots,L_{\mathrm{syn}}$ indexes the resampled
phase grid, distinct from the physical frame grid $0:T$.
These named channels are an action description, not 48 independent physical
causes: some channels are deterministic functions of others, and contact and
phase channels include discrete values. The ambient space $\R^{48}$ describes
numeric storage; neither a nonsingular continuous density nor a Markov
property for $\mathbf d_t$ follows from the simulator's state dynamics.

\paragraph{Frozen video observations.}
For the real-video action learner, the observed variables are features of
an action-providing video $V_{\mathrm r}$, which can be the paired real target.
Let $H_{\ell,\sigma}(V_{\mathrm r};\mathbf c_{\mathrm{ext}},\boldsymbol\xi_{\mathrm{ext}})$ be frozen backbone
features at layer $\ell$ and noise level $\sigma$, with extraction conditioning
$\mathbf c_{\mathrm{ext}}$ and noise $\boldsymbol\xi_{\mathrm{ext}}$. Holding these extraction
settings fixed, the implemented observation operators are
\begin{equation}
 \begin{aligned}
  \mathbf O^{(v)}_{1:L}(V_{\mathrm r})
   &=\mathcal M_v\!\left(
       \bigl(H_{\ell,\sigma}(V_{\mathrm r};\mathbf c_{\mathrm{ext}},\boldsymbol\xi_{\mathrm{ext}})\bigr)_{\sigma\in\Sigma}
     \right),\\
  &\hspace{12mm}v\in\{1,2,\mathrm{full}\},\qquad
       \mathbf O^{(v)}_{1:L}(V_{\mathrm r})\in\R^{L\times d_o}.
 \end{aligned}
 \label{eq:impl-feature-observation}
\end{equation}
Here $\Sigma$ is the fixed ordered list of three extraction noise levels,
and $\mathcal M_v$ includes projection, spatial aggregation, temporal
centering, lag statistics, and fixed training-set standardization.
Views $1,2$ use alternating spatial patches of the same video during
representation training. The full observation uses all spatial patches
and is used to construct the controller packet.
For this real-video model, $L=L_{\mathrm{real}}=21$ and $d_o=3072$.
In the representation equations below, $t$ indexes these resampled phases;
$\sigma$ denotes diffusion noise level and $k$ a sampler step.

The observation operator retains centered global and local statistics,
signed global changes, and local change magnitudes at lags $1,2,4$.
It does not construct an action signal by subtracting a different video.
Both training views inherit information from the same backbone and scene;
they are not assumed conditionally independent given action. In particular,
temporal centering removes a constant feature component but need not
remove content-dependent motion, camera motion, or appearance--motion
interactions.

\subsection{Implemented temporal representation models}
\label{subsec:impl-representation-model}

\paragraph{Real-video action model.}
The Action-VAE is a sequential variational autoencoder~\citep{kingma2014vae}.
Its encoder is a bidirectional GRU~\citep{cho2014gru} with diagonal Gaussian outputs.
For an observed sequence $\mathbf O\in\R^{L\times d_o}$, its posterior density is
\begin{equation}
  q_\phi(\widetilde{\mathbf z}^{\mathrm a}_{1:L}\mid \mathbf O)
  =\prod_{t=1}^{L}
    \mathcal N\!\left(\widetilde{\mathbf z}^{\mathrm a}_t;\boldsymbol\mu_{\phi,t}(\mathbf O),
       \operatorname{diag}(\exp \mathbf s_{\phi,t}(\mathbf O))\right),
  \qquad \widetilde{\mathbf z}^{\mathrm a}_t\in\R^{d_z},\quad d_z=32.
  \label{eq:impl-posterior}
\end{equation}
Here $\boldsymbol\mu_{\phi,t}(\mathbf O)$ and $\mathbf s_{\phi,t}(\mathbf O)$ are $d_z$-dimensional mean and
log-variance vectors, including the implemented log-variance clipping;
$\exp$ acts elementwise. Each posterior mean can depend on the entire
observed sequence. Distinguish a sample $\widetilde{\mathbf z}^{\mathrm a}\sim q_\phi(\cdot\mid \mathbf O)$ from
the deterministic mean code $\widehat{\mathbf z}^{\mathrm a}=\boldsymbol\mu_\phi(\mathbf O)$.
The prior has a recurrent state $\mathbf r_t\in\R^{d_z}$ and conditional densities
\begin{equation}
 \begin{aligned}
  \mathbf r_1&=\mathbf 0,\qquad
  \mathbf r_t=\operatorname{GRU}_\psi(\widetilde{\mathbf z}^{\mathrm a}_{t-1},\mathbf r_{t-1}),\quad t\geq2,\\
  p_{\psi,t}(\widetilde{\mathbf z}^{\mathrm a}_t\mid \widetilde{\mathbf z}^{\mathrm a}_{<t})
    &=\mathcal N\!\left(\widetilde{\mathbf z}^{\mathrm a}_t;\mathbf r_t,\operatorname{diag}(\exp \mathbf s_\psi)\right),\\
  p_\psi(\widetilde{\mathbf z}^{\mathrm a}_{1:L})&=\prod_{t=1}^{L}p_{\psi,t}(\widetilde{\mathbf z}^{\mathrm a}_t\mid \widetilde{\mathbf z}^{\mathrm a}_{<t}).
 \end{aligned}
 \label{eq:impl-recurrent-prior}
\end{equation}
The clipped log-variance vector $\mathbf s_\psi$ is shared across phases.
The mean depends on history through a dense recurrent transition; the product
is a conditional factorization, not temporal independence.
A separate unidirectional GRU decoder $d_\omega$ predicts observations from
a code sequence, with its uncentered output at phase $t$ depending on
$\widetilde{\mathbf z}^{\mathrm a}_{\leq t}$. Neither the prior nor the decoder is a sparse,
text-conditioned structural equation.

The active intrinsic-action training objective evaluates the decoder and
prior at posterior means $\widehat{\mathbf z}^{\mathrm a,(v)}=\boldsymbol\mu_\phi(\mathbf O^{(v)})$.
The representation objective in Eq.~\ref{eq:method-anchor-loss} groups
its auxiliary terms as
\begin{equation}
 \mathcal L_{\mathrm{reg}}
 =\beta\mathcal L_{\mathrm{prior}}
  +\lambda_{\mathrm{ord}}\mathcal L_{\mathrm{ord}}
  +\lambda_{\mathrm{stat}}\mathcal L_{\mathrm{stat}}
  +\mathcal R_{\mathrm{temp}}.
 \label{eq:impl-anchor-regularization}
\end{equation}
Here $\beta$ warms up the prior penalty; the loss components are defined below.
For one two-view example, the reconstruction and prior penalties are
\begin{equation}
 \begin{aligned}
  \ell_{\mathrm{obs}}
    &=\tfrac12\sum_{v=1}^{2}
       \rho\!\left(\operatorname{Center}(d_\omega(\widehat{\mathbf z}^{\mathrm a,(v)}))
          -\operatorname{Center}\!\left(\tfrac{\mathbf O^{(1)}+\mathbf O^{(2)}}2\right)\right),\\
  \ell_{\mathrm{prior}}
    &=\frac{1}{2Ld_z}\sum_{v=1}^{2}\sum_{t=1}^{L}
      \KL\!\left(q_{\phi,t}(\cdot\mid \mathbf O^{(v)})
        \,\|\,p_{\psi,t}(\cdot\mid\widehat{\mathbf z}^{\mathrm a,(v)}_{<t})\right).
 \end{aligned}
 \label{eq:impl-mean-prior-loss}
\end{equation}
For an array $\mathbf W$ with rows $\mathbf w_t$, $\operatorname{Center}(\mathbf W)_t=\mathbf w_t-L^{-1}\sum_{j=1}^L\mathbf w_j$;
$\rho$ averages the elementwise SmoothL1 penalty.
The standard multivariate $\KL$ sums latent-coordinate contributions;
the explicit $1/d_z$ matches their mean in the implementation.
Training averages these per-example penalties over the batch to obtain
the corresponding $\mathcal L$ terms in
Eqs.~\ref{eq:method-anchor-loss} and~\ref{eq:impl-anchor-regularization}.

\paragraph{Auxiliary representation losses.}
Write $\check{\mathbf z}^{\mathrm a,(v)}=\operatorname{Center}(\widehat{\mathbf z}^{\mathrm a,(v)})$ and
$\Delta_\ell\mathbf w_t=\mathbf w_{t+\ell}-\mathbf w_t$ for valid phases. The action signature is
the concatenation of the vectorized centered code and its lag-$1,2,4$
differences, normalized to unit norm with a numerical floor of $10^{-8}$.
Let $s_+$ be the cosine similarity of the two forward-view signatures,
and $s_{\mathrm{rev}},s_{\mathrm{shuf}}$ compare the first forward signature
with signatures obtained by encoding reversed and shuffled observations.
The per-example view and order losses are
\begin{align}
 \ell_{\mathrm{view}}
   &=\rho(\check{\mathbf z}^{\mathrm a,(1)}-\check{\mathbf z}^{\mathrm a,(2)})
     +\tfrac14\rho(\Delta_1\check{\mathbf z}^{\mathrm a,(1)}-\Delta_1\check{\mathbf z}^{\mathrm a,(2)})
     +[m_{\mathrm{view}}-s_+]_+,\nonumber\\
 \ell_{\mathrm{ord}}
   &=\tfrac12\sum_{j\in\{\mathrm{rev},\mathrm{shuf}\}}
       [m_{\mathrm{ord}}-s_++s_j]_+,
 \label{eq:impl-view-order-loss}
\end{align}
where $[x]_+=\max(x,0)$. Shuffling swaps the first two seven-phase blocks
of the extracted trajectory. Static controls repeat either the first or
middle video frame and produce both patch views. Their mean code energy
$e_{\mathrm{stat}}$ and the mean forward-view energy $e_{\mathrm{dyn}}$
use $e(\mathbf W)=\sqrt{(Ld_z)^{-1}\|\operatorname{Center}(\mathbf W)\|_F^2+10^{-8}}$.
The static penalty is
$\ell_{\mathrm{stat}}=e_{\mathrm{stat}}+
[m_{\mathrm{stat}}+e_{\mathrm{stat}}-e_{\mathrm{dyn}}]_+$.

The temporal-difference penalty averages
$\sqrt{(\Delta_\ell\widehat z^{\mathrm a,(1)}_{t,j})^2+10^{-6}}/\ell$
over valid phases, coordinates, and lags $\ell\in\{1,2,4\}$.
The variance penalty averages the deficit below a per-coordinate temporal
standard-deviation floor. The covariance penalty divides each centered coordinate by its temporal
standard deviation, floored at $10^{-4}$, and forms the within-video
covariance with normalization $1/L$. After setting its diagonal to zero,
it averages the squared entries over all $d_z^2$ positions.
Both variance and covariance penalties are averaged over the two views.
The covariance term discourages redundant temporal channels; it neither
requires exact decorrelation nor implies statistical independence.
For $L=21$ and $d_z=32$, a time-centered covariance $\mathbf C$ has rank
at most $20$. If every channel has temporal standard deviation at least
$10^{-4}$, the normalization gives $C_{jj}=1$, and
\[
 \frac{1}{d_z^2}\sum_{i\ne j}C_{ij}^2
 =\frac{\|\mathbf C\|_F^2-d_z}{d_z^2}
 \geq \frac{1}{L-1}-\frac{1}{d_z}=0.01875.
\]
The inequality follows from
$\|\mathbf C\|_F^2\geq(\operatorname{tr}\mathbf C)^2/20$.
Channels affected by the standard-deviation floor need not have unit
variance, so this numerical lower bound is conditional on the stated
normalization regime. The weighted variance, covariance, and difference
penalties form $\mathcal R_{\mathrm{temp}}$.
The default representation configuration uses
$(\lambda_{\mathrm{obs}},\lambda_{\mathrm{view}},\lambda_{\mathrm{ord}},
\lambda_{\mathrm{stat}})=(0.15,1,0.25,0.25)$;
$\beta$ increases linearly to $0.002$ over 400 updates.
The temporal-difference, variance, and covariance weights are
$0.002$, $0.05$, and $0.005$, respectively. The view, order, and static
margins are $0.10$, $0.20$, and $0.05$, and the standard-deviation floor
is $0.08$.
The term named sparse transition in the implementation is this
finite-difference penalty.
Although the architecture supports latent sampling, this training routine
uses posterior means throughout these reconstruction and prior terms.
Consequently, the implemented objective is a regularized representation
loss, not an exact Monte Carlo ELBO for Eq.~\ref{eq:impl-recurrent-prior}.

\paragraph{Supervised synthetic action model.}
The synthetic learner is a separate model with latent
$\widetilde{\mathbf z}^{\mathrm a,\mathrm{syn}}_t\in\R^{48}$, $t=1,\ldots,L_{\mathrm{syn}}$.
Its posterior mean $\widehat{\mathbf z}^{\mathrm a,\mathrm{syn}}$ is directly supervised by $\mathbf d_{1:L_{\mathrm{syn}}}$ in
Eq.~\ref{eq:impl-synthetic-descriptor}.
Its lag-2 prior predicts Gaussian means and variances from the previous
two latent vectors. For $t\geq3$, an alternative adds a strictly
lower-triangular coefficient matrix $\mathbf B$:
\begin{equation}
 \begin{aligned}
  \widetilde z^{\mathrm a,\mathrm{syn}}_{t,j}
   &=m^{\mathrm{syn}}_{\psi,j}(\widetilde{\mathbf z}^{\mathrm a,\mathrm{syn}}_{t-2},\widetilde{\mathbf z}^{\mathrm a,\mathrm{syn}}_{t-1})
     +\tanh\!\left(\sum_{i<j}B_{ji}\widetilde z^{\mathrm a,\mathrm{syn}}_{t,i}\right)\\
   &\quad+\sigma^{\mathrm{syn}}_{\psi,j}
       (\widetilde{\mathbf z}^{\mathrm a,\mathrm{syn}}_{t-2},\widetilde{\mathbf z}^{\mathrm a,\mathrm{syn}}_{t-1})\,\eta^{\mathrm{syn}}_{t,j}.
 \end{aligned}
 \label{eq:impl-synthetic-prior}
\end{equation}
The scales
$\sigma^{\mathrm{syn}}_{\psi,j}$ are positive standard deviations, and the
$\eta^{\mathrm{syn}}_{t,j}$ are independent standard Gaussian noises.
This specifies conditional transitions for $t\geq3$, not an initial law
for the first two phases. It gives a sequential sampling interpretation
of the triangular transition. In training, prior parameters are evaluated
using posterior means, and the instantaneous Jacobian receives an L1 penalty.
This is a model for a learned code, not the simulator's true physical law.
The supervised loss fixes semantic coordinate names; recovery in this
experiment must therefore be distinguished from unsupervised
identifiability. Gaussian priors on the learned code do not imply that
the mixed discrete and continuous descriptor has Gaussian innovations.

\subsection{Connection to editing and the identification question}
\label{subsec:impl-editing-evidence}

The real-video controller consumes a detached packet extracted from the
full observation. Write $\widehat{\mathbf z}^{\mathrm a}(V_{\mathrm r})=\boldsymbol\mu_\phi(\mathbf O^{(\mathrm{full})}(V_{\mathrm r}))$
for its mean code, arranged as an $L\times32$ array at inference, with $L=21$.
The packet is
\begin{equation}
 \mathbf K=\operatorname{stopgrad}\!\left(
   \operatorname{Concat}_{\mathrm{feat}}\!\left(
    \operatorname{Norm}_{\mathrm{phase}}(\widehat{\mathbf z}^{\mathrm a}(V_{\mathrm r})),
    \Psi_{\mathrm J}(V_{\mathrm r})\right)\right)
 \in\R^{L\times64}.
 \label{eq:impl-action-packet}
\end{equation}
Here $\Psi_{\mathrm J}(V_{\mathrm r})\in\R^{L\times32}$ is a separately projected,
ordered V-JEPA2 signal. The two arrays are concatenated along the feature
axis; they are not constrained to be the same identifiable coordinates.
For a noncollapsed $\mathbf W\in\R^{L\times d}$,
$\operatorname{Norm}_{\mathrm{phase}}(\mathbf W)=\operatorname{Center}(\mathbf W)/
\operatorname{RMS}(\operatorname{Center}(\mathbf W))$, where
$\operatorname{RMS}(\mathbf W)=\|\mathbf W\|_F/\sqrt{Ld}$.
The implementation rejects a centered RMS at or below $10^{-8}$.

The source supplies both its video latent and a spatial context
$b_{\mathrm{src}}(V_{\mathrm s})$ formed by interpolating four independently encoded
source frames. This context concatenates absolute features with their
temporal deviations, so it contains source motion as well as appearance.
The resulting editor is conditioned as
$\widehat V\sim p_\theta(\cdot\mid V_{\mathrm s},i,\mathbf K)$.
The packet modulates low-rank updates to the generated branch's
visual self-attention value and output projections, together with
the source spatial context. The analyzed checkpoint injects these
updates during the first eight steps of a native 40-step sampler.

The controller uses a spatial warm-up followed by an early-step behavior
stage. Warm-up retains Action-VAE output alignment and a periodic
V-JEPA2 critic. The subsequent stage retains the frozen Action-VAE
condition but disables its output-matching loss. Its training queries
come from native sampler prefixes. Distinguish a noisy video latent
$\boldsymbol\ell_k\in\mathcal H$ from a decoded video
$\widehat V_{0\mid k}\in\mathcal V$:
\begin{equation}
 \begin{aligned}
  \boldsymbol\ell_k&=\operatorname{Prefix}_k(V_{\mathrm s},i,\mathbf K;\boldsymbol\xi),\\
  \widehat{\boldsymbol\ell}_{0\mid k}
    &=\boldsymbol\ell_k-\sigma_k v_\theta^{\mathrm{guided}}
          (\boldsymbol\ell_k,\sigma_k;V_{\mathrm s},i,\mathbf K,b_{\mathrm{src}}(V_{\mathrm s})),\\
  \widehat V_{0\mid k}
    &=\operatorname{Decode}_{\mathrm{VAE}}(\widehat{\boldsymbol\ell}_{0\mid k}).
 \end{aligned}
 \label{eq:impl-controller-query}
\end{equation}
The clean-latent relation uses the velocity convention
$\boldsymbol\ell_\sigma=(1-\sigma)\boldsymbol\ell_0+\sigma\boldsymbol\epsilon_{\mathrm{diff}}$,
$\mathbf v=\boldsymbol\epsilon_{\mathrm{diff}}-\boldsymbol\ell_0$.
The effective guided velocity includes the native conditional and
unconditional branches. The decoding operator includes the fixed query
view selection, latent denormalization, and video-VAE decoding used by
the observer. Training cycles through prefix indices $k\in\{0,2,4,6\}$
and evaluates these decoded clean predictions; it does not differentiate
through a complete 40-step rollout.

Let $e_{\mathrm J}$ be the frozen ordered V-JEPA2 behavior reader.
For each query, index the comparison videos by $\mathcal K=\{+,\mathrm{src},\mathrm{rev},\mathrm{shuf}\}$,
with $V_+=V_{\mathrm r}$, $V_{\mathrm{src}}=V_{\mathrm s}$, $V_{\mathrm{rev}}=\operatorname{Reverse}(V_{\mathrm r})$,
and $V_{\mathrm{shuf}}=\operatorname{Shuffle}(V_{\mathrm r})$.
Define
$\gamma_{j,k}=\operatorname{sim}(e_{\mathrm J}(\widehat V_{0\mid k}),e_{\mathrm J}(V_j))$,
where $\operatorname{sim}$ is the implemented mean cosine comparison
between corresponding descriptor vectors.
The primary behavior penalty is
\begin{equation}
 \ell_{\mathrm{act},k}
 =1-\gamma_{+,k}-\lambda_{\mathrm{disc}}\log
   \frac{\exp(\gamma_{+,k}/\tau)}
        {\sum_{j\in\mathcal K}\exp(\gamma_{j,k}/\tau)},
 \qquad \tau=\lambda_{\mathrm{disc}}=0.10.
 \label{eq:impl-behavior-loss}
\end{equation}
The candidate indices preserve all four terms even if two candidate
videos coincide. The temperature $\tau$ and discrimination weight
$\lambda_{\mathrm{disc}}$ have distinct roles despite their equal values.
The training objective $\mathcal L_{\mathrm{act}}$ averages
$\ell_{\mathrm{act},k}$ over training examples and the prescribed query cycle.
Reference and control descriptors are detached; gradients through the
generated-video descriptor remain active for all four comparisons.
This high-dimensional behavior reader differs from both the
Action-VAE code and the compressed V-JEPA2 packet component.

\paragraph{Preservation-gradient combination.}
Source preservation uses symmetric patch-set comparisons and bounded
appearance and temporal statistics. Let $\mathbf g_{\mathrm a}$ and
$\mathbf g_{\mathrm p}$ denote the behavior and preservation gradients with
respect to the clean-latent prediction
$\widehat{\boldsymbol\ell}_{0\mid k}$, vectorized for the inner product.
The implemented combination is
\begin{equation}
 \begin{aligned}
  \widetilde{\mathbf g}_{\mathrm p}
    &=\mathbf g_{\mathrm p}
      -\frac{\min\{\langle\mathbf g_{\mathrm a},\mathbf g_{\mathrm p}\rangle,0\}}
             {\max\{\|\mathbf g_{\mathrm a}\|^2,\varepsilon_g\}}
       \mathbf g_{\mathrm a},\\
  s&=\min\!\left\{1,
    \frac{\eta\,\operatorname{RMS}(\mathbf g_{\mathrm a})}
         {\max\{\operatorname{RMS}(\widetilde{\mathbf g}_{\mathrm p}),\varepsilon_g\}}
       \right\},\qquad
  \mathbf g=\mathbf g_{\mathrm a}+s\widetilde{\mathbf g}_{\mathrm p},
 \end{aligned}
 \label{eq:impl-gradient-combination}
\end{equation}
where $\eta=1/4$ and $\varepsilon_g=10^{-24}$ is the numerical floor.
The cap never amplifies the preservation gradient. Projection acts in
clean-prediction space. Backpropagation maps each gradient to
$\mathbf J^\top\mathbf g_{\bullet}$, where
$\mathbf J=\partial\widehat{\boldsymbol\ell}_{0\mid k}/\partial\theta$.
The resulting parameter-space inner product uses
$\mathbf J\mathbf J^\top$, so nonnegative prediction-space inner products
need not remain nonnegative.
This rule therefore does not guarantee monotonic decrease of the behavior
loss under parameter updates. It is not a fixed weighted sum of losses.
The stage does not use a paired target as a clean flow-matching endpoint.

Inference uses the ordinary conditional sampler. The source of its
reference packet is specified in Appendix~\ref{subsec:impl-configuration}.

\paragraph{Identification target.}
The identification criterion concerns the deterministic mean code.
Let $\mathbf z^{\mathrm a}_{\mathrm r,[L]}=\mathcal R_L(\mathbf z^{\mathrm a}_{\mathrm r})$ denote the task-relevant behavior on the
representation grid. When this resampling retains the behavioral information
of interest, the corresponding block identification criterion is
\begin{equation}
  \widehat{\mathbf z}^{\mathrm a}(V_{\mathrm r})=h_{\mathrm a,L}(\mathbf z^{\mathrm a}_{\mathrm r,[L]}),
  \label{eq:impl-identification-target}
\end{equation}
where $h_{\mathrm a,L}$ is shared across compatible content contexts,
injective on the supported resampled action space, and invertible on its
image. Equation~\ref{eq:impl-identification-target} defines the criterion;
its attainment depends on the observation and view-learning conditions.
The normalized controller packet is a different object: its centering
and RMS normalization discard offsets and overall positive scale.

Appendix~\ref{subsec:conditioning-proof} analyzes the implemented views,
prediction target, and complete packet against this identification
criterion. The evaluation protocol in Appendix~\ref{subsec:synthetic-protocol-details}
measures action retention and content dependence before assessing final
editing behavior. The source context also contains motion and is not an
identified content code; the two-code decoder in
Appendix~\ref{sec:controlled-transfer} is a separate controlled baseline.
The factorized recurrent prior supplies temporal regularization.
Coordinate-level identification would additionally require compatible
observation and transition models and sufficient distributional
variation~\citep{yao2022tdrl,shen2026covogan}.

\subsection{Reproducible configuration and input settings}
\label{subsec:impl-configuration}

Table~\ref{tab:impl-configuration} specifies the implementation. The
editor is trained on 1,508 source--target pairs, about 16\% of the pairs in
the released real-world dataset
(Appendix~\ref{subsec:real-dataset}). Training demonstrations mix the
real target clips of these pairs with reference videos of the target
actions generated offline, in a 3:1 ratio of real to generated videos.
Training with a 1:1 ratio gave similar qualitative results and
OmniEdit-Bench Temporal scores, so within this range the editor is not
sensitive to the proportion of real targets; we use 3:1 for all reported
results. Demonstrations are read only by the frozen extractors that
produce the action condition and the behavior descriptors.

\begin{table}[!htbp]
 \centering
 \small
 \renewcommand{\arraystretch}{1.15}
 \caption{Implementation configuration.}
 \label{tab:impl-configuration}
 \begin{tabular}{@{}p{0.23\linewidth}p{0.73\linewidth}@{}}
 \hline
 Component & Configuration \\ \hline
 Frozen editor & Bernini-R-1.3B (Diffusers release); VeOmni v0.1.11. \\
 Feature extraction & Block index 24; noise levels $(0.85,0.55,0.20)$;
 frozen 128-dimensional projections; 81 sampled frames,
 $L=21$ temporal phases and $d_o=3072$. \\
 Action-VAE & Two-layer bidirectional GRU, 128 hidden units per direction;
 two-layer decoder GRU with 256 hidden units; dropout 0.10;
 32-dimensional posterior and recurrent prior state. \\
 Representation training & AdamW~\citep{loshchilov2019adamw}, 1,200 updates, batch 8,
 learning rate $3\!\times\!10^{-4}$, weight decay $10^{-4}$,
 gradient-norm cap 5. \\
 Behavior reader & Frozen V-JEPA2 ViT-L (64-frame clips at 256 pixels);
 last-layer ordered descriptors. \\
 Reference condition & $21\!\times\!32$ Action-VAE code concatenated with
 $21\!\times\!32$ compressed V-JEPA2 features. \\
 Adapters and source & Rank 64, scale parameter 64; visual self-attention
 value/output projections in blocks 0--24. Source-frame indices
 $(0,27,53,80)$ provide spatial context. \\
 Controller training & 1,508 training pairs; real targets and offline-generated
 references as demonstrations at 3:1; a warm-up stage followed by
 behavior supervision. AdamW learning rates $5\!\times\!10^{-4}$ and
 $2\!\times\!10^{-4}$, respectively; weight decay 0, gradient-norm cap 1. \\
 Sampling & Native 40-step sampler; action adapters active at indices
 $0\leq k<8$. Behavior-training queries use $k\in\{0,2,4,6\}$. \\
 Inference demonstration & Qwen3-VL rewrites the source and instruction as a
 text-to-video prompt; the frozen editor generates the demonstration from
 this prompt; no real target is read. \\
 \hline
 \end{tabular}
\end{table}

\paragraph{Feature and packet construction.}
Extraction samples 81 frames uniformly over the full clip and uses a
maximum pixel budget of 245,760 with stride 16, presented to the model
at 16 fps. The $1536\!\times\!128$ feature projection has independent
random-sign entries scaled by $1/\sqrt{128}$ and a fixed seed.
The noise realization is fixed for each example and instruction and
reused across the three noise levels. Extraction uses the recorded editing instruction.
The V-JEPA2 branch resamples its $32\!\times\!1024$ ordered sequence to
21 phases by linear interpolation with aligned endpoints. The level
sequence and lag-$1,2,4$ differences each use a fixed
$1024\!\times\!8$ random-sign projection, scaled by $1/\sqrt8$, with
its own fixed seed. Lag differences are zero-padded
at their unavailable initial phases. Their concatenation is centered
and RMS-normalized to form $\Psi_{\mathrm J}$.
Both 32-channel branches therefore have zero temporal mean and unit RMS
before concatenation. Re-normalizing the concatenated packet at injection
preserves these properties up to storage and arithmetic precision.

\paragraph{Controller schedule.}
The warm-up uses native query index 6 and a periodic behavior-reader
query every fourth update. The behavior stage uses the four query indices
in Table~\ref{tab:impl-configuration} and queries the reader on every
example. It inherits the warm-up adapter weights and resets the AdamW
state.

\paragraph{Inference with generated demonstrations.}
At inference, no real target is available.
Qwen3-VL~\citep{bai2025qwen3vl} first rewrites the source video and
editing instruction as a text-to-video prompt describing the requested
action, and the frozen editor, Bernini-R-1.3B, generates a demonstration
from this prompt in text-to-video mode.
The same frozen extractors compute the action condition $\mathbf K$ from
this generated video, and the trained adapters then edit the source with
$\mathbf K$ using the native 40-step sampler, active only in the first
eight steps. Neither benchmark targets nor other real targets are read at
inference, and no external video generator is used; obtaining the demonstration
adds one vision--language query and one sampling pass of the frozen
editor.

The planned intervention tests will use independent appearance variants
and separately trained Action-VAE-only, V-JEPA2-only, and combined editors
(Appendix~\ref{subsec:synthetic-protocol-details}). Final videos will be
assessed for action execution and source fidelity using independent
behavior measurements and human review; intermediate reader agreement
is a separate diagnostic. Benchmark comparisons are reported in
Section~\ref{subsec:real-experiments}; this human assessment is not yet
complete in the present draft.

